\begin{figure}[t]
\centering
\resizebox{\linewidth}{!}{%
\begin{tikzpicture}[
  font=\small,
  panel/.style={draw=black!18, rounded corners=3pt, fill=black!2},
  chip/.style={draw=black!25, rounded corners=2pt, fill=white, align=center, inner xsep=5pt, inner ysep=3pt},
  good/.style={draw=green!55!black, fill=green!12},
  bad/.style={draw=orange!70!black, fill=orange!12},
  arrow/.style={-{Latex[length=2mm]}, line width=0.75pt},
  lab/.style={font=\footnotesize, text=black!70}
]
\node[panel, minimum width=6.9cm, minimum height=4.2cm, anchor=south west] at (0,0) {};
\node[anchor=west, font=\bfseries] at (0.25,3.82) {救回证据：p-crAssphage};
\node[chip] (h54) at (1.15,2.55) {Hybrid\\第 54 位};
\node[chip, good] (k1) at (5.75,2.55) {KCH-MedRank\\第 1 位};
\draw[arrow, draw=green!50!black] (h54) -- node[above, lab] {提升至 top 10} (k1);
\node[chip, good] at (1.25,1.25) {bacteriophage};
\node[chip, good] at (3.15,1.25) {crAssphage};
\node[chip, good] at (5.0,1.25) {viral genome};
\node[lab, align=center] at (3.45,0.48) {共享生物医学概念和局部 phage 相关候选簇\\共同支持证据提升。};

\node[panel, minimum width=6.9cm, minimum height=4.2cm, anchor=south west] at (7.35,0) {};
\node[anchor=west, font=\bfseries] at (7.6,3.82) {丢失证据：ubiquitination};
\node[chip, bad] (h2) at (8.45,2.55) {Hybrid\\第 2 位};
\node[chip] (k18) at (13.05,2.55) {KCH-MedRank\\第 18 位};
\draw[arrow, draw=orange!75!black] (h2) -- node[above, lab] {降出 top 10} (k18);
\node[chip, bad] at (8.65,1.25) {mono-/poly-\\ubiquitination};
\node[chip, bad] at (10.85,1.25) {ubiquitin};
\node[chip, bad] at (12.85,1.25) {post-translational\\processing};
\node[lab, align=center] at (10.85,0.48) {表面形式不匹配和概念归一化不足\\削弱了知识对齐信号。};
\end{tikzpicture}}
\caption{BioASQ 中一个救回证据和一个丢失证据的可视化案例。该图概括了局部生物医学概念对齐如何提升被遗漏的金标准段落，同时也显示同义词与层级对齐不充分时可能降低相关段落排名。}
\label{fig:case-study-path}
\end{figure}
